\documentclass[11pt]{article}
% jugeo-common.tex — shared preamble for all Judgment Geometry papers
% Include via: % jugeo-common.tex — shared preamble for all Judgment Geometry papers
% Include via: % jugeo-common.tex — shared preamble for all Judgment Geometry papers
% Include via: \input{jugeo-common}

% ── Packages ────────────────────────────────────────────────────────────
\usepackage{amsmath,amssymb,amsthm,mathtools}
\usepackage{tikz-cd}
\usepackage{listings}
\usepackage{xcolor}
\usepackage{xspace}
\usepackage{booktabs}
\usepackage{multirow}
\usepackage{enumitem}
\usepackage[expansion=false]{microtype}
\usepackage{hyperref}
\usepackage{cleveref}
% \usepackage{thmtools}  % disabled: not available in this TeX installation
\usepackage{float}
\usepackage{caption}
\usepackage{subcaption}
\usepackage{tabularx}
\usepackage{stmaryrd}       % \llbracket, \rrbracket
\usepackage{bbm}            % \mathbbm{1}

% ── Page geometry ───────────────────────────────────────────────────────
\usepackage[margin=1in]{geometry}

% ── Theorem environments ────────────────────────────────────────────────
\theoremstyle{plain}
\newtheorem{theorem}{Theorem}[section]
\newtheorem{lemma}[theorem]{Lemma}
\newtheorem{proposition}[theorem]{Proposition}
\newtheorem{corollary}[theorem]{Corollary}

\theoremstyle{definition}
\newtheorem{definition}[theorem]{Definition}
\newtheorem{example}[theorem]{Example}
\newtheorem{construction}[theorem]{Construction}

\theoremstyle{remark}
\newtheorem{remark}[theorem]{Remark}
\newtheorem{notation}[theorem]{Notation}

% ── Python listings style ──────────────────────────────────────────────
\definecolor{jugeoBlue}{HTML}{2563EB}
\definecolor{jugeoGreen}{HTML}{059669}
\definecolor{jugeoRed}{HTML}{DC2626}
\definecolor{jugeoPurple}{HTML}{7C3AED}
\definecolor{jugeoGray}{HTML}{6B7280}
\definecolor{jugeoBg}{HTML}{F8FAFC}

\lstdefinestyle{jugeo-python}{
  language=Python,
  basicstyle=\ttfamily\small,
  keywordstyle=\color{jugeoBlue}\bfseries,
  stringstyle=\color{jugeoGreen},
  commentstyle=\color{jugeoGray}\itshape,
  numberstyle=\tiny\color{jugeoGray},
  backgroundcolor=\color{jugeoBg},
  numbers=left,
  numbersep=6pt,
  frame=single,
  framerule=0.4pt,
  rulecolor=\color{jugeoGray!40},
  breaklines=true,
  breakatwhitespace=true,
  showstringspaces=false,
  tabsize=4,
  xleftmargin=1.5em,
  framexleftmargin=1.5em,
  morekeywords={dataclass,frozen,field,Enum,IntEnum,auto,Any,
                Mapping,Sequence,Optional,match,case},
  literate={→}{{$\to$}}1 {←}{{$\leftarrow$}}1
           {≤}{{$\leq$}}1 {≥}{{$\geq$}}1 {≠}{{$\neq$}}1
           {∈}{{$\in$}}1 {∀}{{$\forall$}}1 {∃}{{$\exists$}}1
           {⊕}{{$\oplus$}}1 {⊖}{{$\ominus$}}1,
  escapeinside={(*@}{@*)},
}
\lstset{style=jugeo-python}

\lstdefinestyle{jugeo-output}{
  basicstyle=\ttfamily\small,
  backgroundcolor=\color{jugeoBg},
  frame=single,
  framerule=0.4pt,
  rulecolor=\color{jugeoGray!40},
  breaklines=true,
  showstringspaces=false,
  xleftmargin=1.5em,
  framexleftmargin=0.5em,
  numbers=none,
}

% ── JuGeo-specific macros ──────────────────────────────────────────────

% Judgment 8-tuple
\newcommand{\J}{\mathcal{J}}
\newcommand{\Jud}[8]{(#1,\,#2,\,#3,\,#4,\,#5,\,#6,\,#7,\,#8)}
\newcommand{\JudShort}[1]{\J_{#1}}

% Components
\newcommand{\coord}{\mathsf{c}}            % coordinate
\newcommand{\prop}{\varphi}                 % proposition
\newcommand{\carrier}{\mathcal{A}}          % carrier type
\newcommand{\evid}{\mathcal{E}}             % evidence bundle
\newcommand{\oblig}{\mathcal{O}}            % obligations
\newcommand{\obstr}{\mathcal{B}}            % obstructions
\newcommand{\trust}{\mathcal{T}}            % trust annotation
\newcommand{\prov}{\Pi}                     % provenance

% Trust algebra
\newcommand{\Talg}{\mathfrak{T}}            % trust ordered algebra
\newcommand{\Eadm}{\mathcal{E}_{\mathrm{adm}}}
\newcommand{\tleq}{\preceq}                 % trust ordering
\newcommand{\tjoin}{\oplus}                 % conservative join
\newcommand{\tatten}{\ominus}               % attenuation
\newcommand{\tpromote}{\uparrow_\pi}        % promotion
\newcommand{\tdemote}{\downarrow_\chi}       % demotion

% Trust tiers
\newcommand{\Tcontra}{\ensuremath{\mathsf{CONTRADICTED}}}
\newcommand{\Tunver}{\ensuremath{\mathsf{UNVERIFIED}}}
\newcommand{\Tcopilot}{\ensuremath{\mathsf{COPILOT}}}
\newcommand{\Toracle}{\ensuremath{\mathsf{ORACLE}}}
\newcommand{\Truntime}{\ensuremath{\mathsf{RUNTIME}}}
\newcommand{\Tsolver}{\ensuremath{\mathsf{SOLVER}}}
\newcommand{\Tproof}{\ensuremath{\mathsf{PROOF}}}

% Site and geometry
\newcommand{\Site}{\mathcal{S}}             % semantic site
\newcommand{\Cov}{\mathrm{Cov}}            % covering family
\newcommand{\Ob}{\mathrm{Ob}}              % objects
\newcommand{\Mor}{\mathrm{Mor}}            % morphisms
\newcommand{\res}{\mathrm{res}}            % restriction
\newcommand{\Cech}{\check{C}}              % Čech complex
\newcommand{\Hcoh}[1]{H^{#1}}             % cohomology group

% Descent
\newcommand{\descent}{\mathrm{desc}}
\newcommand{\glue}{\mathrm{glue}}
\newcommand{\overlap}[2]{#1 \cap #2}

% Semantic moves
\newcommand{\VERIFY}{\textsc{Verify}}
\newcommand{\CONSTRUCT}{\textsc{Construct}}
\newcommand{\NORMALIZE}{\textsc{Normalize}}
\newcommand{\GLUE}{\textsc{Glue}}
\newcommand{\RESTRICT}{\textsc{Restrict}}
\newcommand{\TRANSPORT}{\textsc{Transport}}
\newcommand{\REPAIR}{\textsc{Repair}}
\newcommand{\PROMOTE}{\textsc{Promote}}
\newcommand{\CHALLENGE}{\textsc{Challenge}}
\newcommand{\ESCALATE}{\textsc{Escalate}}

% Fragments
\newcommand{\QFLIA}{\texttt{QF\_LIA}}
\newcommand{\QFLRA}{\texttt{QF\_LRA}}
\newcommand{\QFBV}{\texttt{QF\_BV}}
\newcommand{\QFUF}{\texttt{QF\_UF}}

% Misc
\newcommand{\Python}{\textsc{Python}}
\newcommand{\JuGeo}{\textsc{JuGeo}}
\newcommand{\LEAN}{\textsc{Lean}\,4}
\newcommand{\Fstar}{F$^{\star}$}
\newcommand{\Coq}{\textsc{Coq}}
\newcommand{\Dafny}{\textsc{Dafny}}

% Common shortcuts
\newcommand{\ie}{i.e.\@\xspace}
\newcommand{\eg}{e.g.\@\xspace}
\newcommand{\cf}{cf.\@\xspace}
\newcommand{\wrt}{w.r.t.\@\xspace}

% Evaluation shortcuts
\newcommand{\numcases}{300}
\newcommand{\numspec}{100}
\newcommand{\numeq}{100}
\newcommand{\numbug}{100}
\newcommand{\medtime}{6.5\,ms}
\newcommand{\totaltime}{1.949\,s}


% ── Packages ────────────────────────────────────────────────────────────
\usepackage{amsmath,amssymb,amsthm,mathtools}
\usepackage{tikz-cd}
\usepackage{listings}
\usepackage{xcolor}
\usepackage{xspace}
\usepackage{booktabs}
\usepackage{multirow}
\usepackage{enumitem}
\usepackage[expansion=false]{microtype}
\usepackage{hyperref}
\usepackage{cleveref}
% \usepackage{thmtools}  % disabled: not available in this TeX installation
\usepackage{float}
\usepackage{caption}
\usepackage{subcaption}
\usepackage{tabularx}
\usepackage{stmaryrd}       % \llbracket, \rrbracket
\usepackage{bbm}            % \mathbbm{1}

% ── Page geometry ───────────────────────────────────────────────────────
\usepackage[margin=1in]{geometry}

% ── Theorem environments ────────────────────────────────────────────────
\theoremstyle{plain}
\newtheorem{theorem}{Theorem}[section]
\newtheorem{lemma}[theorem]{Lemma}
\newtheorem{proposition}[theorem]{Proposition}
\newtheorem{corollary}[theorem]{Corollary}

\theoremstyle{definition}
\newtheorem{definition}[theorem]{Definition}
\newtheorem{example}[theorem]{Example}
\newtheorem{construction}[theorem]{Construction}

\theoremstyle{remark}
\newtheorem{remark}[theorem]{Remark}
\newtheorem{notation}[theorem]{Notation}

% ── Python listings style ──────────────────────────────────────────────
\definecolor{jugeoBlue}{HTML}{2563EB}
\definecolor{jugeoGreen}{HTML}{059669}
\definecolor{jugeoRed}{HTML}{DC2626}
\definecolor{jugeoPurple}{HTML}{7C3AED}
\definecolor{jugeoGray}{HTML}{6B7280}
\definecolor{jugeoBg}{HTML}{F8FAFC}

\lstdefinestyle{jugeo-python}{
  language=Python,
  basicstyle=\ttfamily\small,
  keywordstyle=\color{jugeoBlue}\bfseries,
  stringstyle=\color{jugeoGreen},
  commentstyle=\color{jugeoGray}\itshape,
  numberstyle=\tiny\color{jugeoGray},
  backgroundcolor=\color{jugeoBg},
  numbers=left,
  numbersep=6pt,
  frame=single,
  framerule=0.4pt,
  rulecolor=\color{jugeoGray!40},
  breaklines=true,
  breakatwhitespace=true,
  showstringspaces=false,
  tabsize=4,
  xleftmargin=1.5em,
  framexleftmargin=1.5em,
  morekeywords={dataclass,frozen,field,Enum,IntEnum,auto,Any,
                Mapping,Sequence,Optional,match,case},
  literate={→}{{$\to$}}1 {←}{{$\leftarrow$}}1
           {≤}{{$\leq$}}1 {≥}{{$\geq$}}1 {≠}{{$\neq$}}1
           {∈}{{$\in$}}1 {∀}{{$\forall$}}1 {∃}{{$\exists$}}1
           {⊕}{{$\oplus$}}1 {⊖}{{$\ominus$}}1,
  escapeinside={(*@}{@*)},
}
\lstset{style=jugeo-python}

\lstdefinestyle{jugeo-output}{
  basicstyle=\ttfamily\small,
  backgroundcolor=\color{jugeoBg},
  frame=single,
  framerule=0.4pt,
  rulecolor=\color{jugeoGray!40},
  breaklines=true,
  showstringspaces=false,
  xleftmargin=1.5em,
  framexleftmargin=0.5em,
  numbers=none,
}

% ── JuGeo-specific macros ──────────────────────────────────────────────

% Judgment 8-tuple
\newcommand{\J}{\mathcal{J}}
\newcommand{\Jud}[8]{(#1,\,#2,\,#3,\,#4,\,#5,\,#6,\,#7,\,#8)}
\newcommand{\JudShort}[1]{\J_{#1}}

% Components
\newcommand{\coord}{\mathsf{c}}            % coordinate
\newcommand{\prop}{\varphi}                 % proposition
\newcommand{\carrier}{\mathcal{A}}          % carrier type
\newcommand{\evid}{\mathcal{E}}             % evidence bundle
\newcommand{\oblig}{\mathcal{O}}            % obligations
\newcommand{\obstr}{\mathcal{B}}            % obstructions
\newcommand{\trust}{\mathcal{T}}            % trust annotation
\newcommand{\prov}{\Pi}                     % provenance

% Trust algebra
\newcommand{\Talg}{\mathfrak{T}}            % trust ordered algebra
\newcommand{\Eadm}{\mathcal{E}_{\mathrm{adm}}}
\newcommand{\tleq}{\preceq}                 % trust ordering
\newcommand{\tjoin}{\oplus}                 % conservative join
\newcommand{\tatten}{\ominus}               % attenuation
\newcommand{\tpromote}{\uparrow_\pi}        % promotion
\newcommand{\tdemote}{\downarrow_\chi}       % demotion

% Trust tiers
\newcommand{\Tcontra}{\ensuremath{\mathsf{CONTRADICTED}}}
\newcommand{\Tunver}{\ensuremath{\mathsf{UNVERIFIED}}}
\newcommand{\Tcopilot}{\ensuremath{\mathsf{COPILOT}}}
\newcommand{\Toracle}{\ensuremath{\mathsf{ORACLE}}}
\newcommand{\Truntime}{\ensuremath{\mathsf{RUNTIME}}}
\newcommand{\Tsolver}{\ensuremath{\mathsf{SOLVER}}}
\newcommand{\Tproof}{\ensuremath{\mathsf{PROOF}}}

% Site and geometry
\newcommand{\Site}{\mathcal{S}}             % semantic site
\newcommand{\Cov}{\mathrm{Cov}}            % covering family
\newcommand{\Ob}{\mathrm{Ob}}              % objects
\newcommand{\Mor}{\mathrm{Mor}}            % morphisms
\newcommand{\res}{\mathrm{res}}            % restriction
\newcommand{\Cech}{\check{C}}              % Čech complex
\newcommand{\Hcoh}[1]{H^{#1}}             % cohomology group

% Descent
\newcommand{\descent}{\mathrm{desc}}
\newcommand{\glue}{\mathrm{glue}}
\newcommand{\overlap}[2]{#1 \cap #2}

% Semantic moves
\newcommand{\VERIFY}{\textsc{Verify}}
\newcommand{\CONSTRUCT}{\textsc{Construct}}
\newcommand{\NORMALIZE}{\textsc{Normalize}}
\newcommand{\GLUE}{\textsc{Glue}}
\newcommand{\RESTRICT}{\textsc{Restrict}}
\newcommand{\TRANSPORT}{\textsc{Transport}}
\newcommand{\REPAIR}{\textsc{Repair}}
\newcommand{\PROMOTE}{\textsc{Promote}}
\newcommand{\CHALLENGE}{\textsc{Challenge}}
\newcommand{\ESCALATE}{\textsc{Escalate}}

% Fragments
\newcommand{\QFLIA}{\texttt{QF\_LIA}}
\newcommand{\QFLRA}{\texttt{QF\_LRA}}
\newcommand{\QFBV}{\texttt{QF\_BV}}
\newcommand{\QFUF}{\texttt{QF\_UF}}

% Misc
\newcommand{\Python}{\textsc{Python}}
\newcommand{\JuGeo}{\textsc{JuGeo}}
\newcommand{\LEAN}{\textsc{Lean}\,4}
\newcommand{\Fstar}{F$^{\star}$}
\newcommand{\Coq}{\textsc{Coq}}
\newcommand{\Dafny}{\textsc{Dafny}}

% Common shortcuts
\newcommand{\ie}{i.e.\@\xspace}
\newcommand{\eg}{e.g.\@\xspace}
\newcommand{\cf}{cf.\@\xspace}
\newcommand{\wrt}{w.r.t.\@\xspace}

% Evaluation shortcuts
\newcommand{\numcases}{300}
\newcommand{\numspec}{100}
\newcommand{\numeq}{100}
\newcommand{\numbug}{100}
\newcommand{\medtime}{6.5\,ms}
\newcommand{\totaltime}{1.949\,s}


% ── Packages ────────────────────────────────────────────────────────────
\usepackage{amsmath,amssymb,amsthm,mathtools}
\usepackage{tikz-cd}
\usepackage{listings}
\usepackage{xcolor}
\usepackage{xspace}
\usepackage{booktabs}
\usepackage{multirow}
\usepackage{enumitem}
\usepackage[expansion=false]{microtype}
\usepackage{hyperref}
\usepackage{cleveref}
% \usepackage{thmtools}  % disabled: not available in this TeX installation
\usepackage{float}
\usepackage{caption}
\usepackage{subcaption}
\usepackage{tabularx}
\usepackage{stmaryrd}       % \llbracket, \rrbracket
\usepackage{bbm}            % \mathbbm{1}

% ── Page geometry ───────────────────────────────────────────────────────
\usepackage[margin=1in]{geometry}

% ── Theorem environments ────────────────────────────────────────────────
\theoremstyle{plain}
\newtheorem{theorem}{Theorem}[section]
\newtheorem{lemma}[theorem]{Lemma}
\newtheorem{proposition}[theorem]{Proposition}
\newtheorem{corollary}[theorem]{Corollary}

\theoremstyle{definition}
\newtheorem{definition}[theorem]{Definition}
\newtheorem{example}[theorem]{Example}
\newtheorem{construction}[theorem]{Construction}

\theoremstyle{remark}
\newtheorem{remark}[theorem]{Remark}
\newtheorem{notation}[theorem]{Notation}

% ── Python listings style ──────────────────────────────────────────────
\definecolor{jugeoBlue}{HTML}{2563EB}
\definecolor{jugeoGreen}{HTML}{059669}
\definecolor{jugeoRed}{HTML}{DC2626}
\definecolor{jugeoPurple}{HTML}{7C3AED}
\definecolor{jugeoGray}{HTML}{6B7280}
\definecolor{jugeoBg}{HTML}{F8FAFC}

\lstdefinestyle{jugeo-python}{
  language=Python,
  basicstyle=\ttfamily\small,
  keywordstyle=\color{jugeoBlue}\bfseries,
  stringstyle=\color{jugeoGreen},
  commentstyle=\color{jugeoGray}\itshape,
  numberstyle=\tiny\color{jugeoGray},
  backgroundcolor=\color{jugeoBg},
  numbers=left,
  numbersep=6pt,
  frame=single,
  framerule=0.4pt,
  rulecolor=\color{jugeoGray!40},
  breaklines=true,
  breakatwhitespace=true,
  showstringspaces=false,
  tabsize=4,
  xleftmargin=1.5em,
  framexleftmargin=1.5em,
  morekeywords={dataclass,frozen,field,Enum,IntEnum,auto,Any,
                Mapping,Sequence,Optional,match,case},
  literate={→}{{$\to$}}1 {←}{{$\leftarrow$}}1
           {≤}{{$\leq$}}1 {≥}{{$\geq$}}1 {≠}{{$\neq$}}1
           {∈}{{$\in$}}1 {∀}{{$\forall$}}1 {∃}{{$\exists$}}1
           {⊕}{{$\oplus$}}1 {⊖}{{$\ominus$}}1,
  escapeinside={(*@}{@*)},
}
\lstset{style=jugeo-python}

\lstdefinestyle{jugeo-output}{
  basicstyle=\ttfamily\small,
  backgroundcolor=\color{jugeoBg},
  frame=single,
  framerule=0.4pt,
  rulecolor=\color{jugeoGray!40},
  breaklines=true,
  showstringspaces=false,
  xleftmargin=1.5em,
  framexleftmargin=0.5em,
  numbers=none,
}

% ── JuGeo-specific macros ──────────────────────────────────────────────

% Judgment 8-tuple
\newcommand{\J}{\mathcal{J}}
\newcommand{\Jud}[8]{(#1,\,#2,\,#3,\,#4,\,#5,\,#6,\,#7,\,#8)}
\newcommand{\JudShort}[1]{\J_{#1}}

% Components
\newcommand{\coord}{\mathsf{c}}            % coordinate
\newcommand{\prop}{\varphi}                 % proposition
\newcommand{\carrier}{\mathcal{A}}          % carrier type
\newcommand{\evid}{\mathcal{E}}             % evidence bundle
\newcommand{\oblig}{\mathcal{O}}            % obligations
\newcommand{\obstr}{\mathcal{B}}            % obstructions
\newcommand{\trust}{\mathcal{T}}            % trust annotation
\newcommand{\prov}{\Pi}                     % provenance

% Trust algebra
\newcommand{\Talg}{\mathfrak{T}}            % trust ordered algebra
\newcommand{\Eadm}{\mathcal{E}_{\mathrm{adm}}}
\newcommand{\tleq}{\preceq}                 % trust ordering
\newcommand{\tjoin}{\oplus}                 % conservative join
\newcommand{\tatten}{\ominus}               % attenuation
\newcommand{\tpromote}{\uparrow_\pi}        % promotion
\newcommand{\tdemote}{\downarrow_\chi}       % demotion

% Trust tiers
\newcommand{\Tcontra}{\ensuremath{\mathsf{CONTRADICTED}}}
\newcommand{\Tunver}{\ensuremath{\mathsf{UNVERIFIED}}}
\newcommand{\Tcopilot}{\ensuremath{\mathsf{COPILOT}}}
\newcommand{\Toracle}{\ensuremath{\mathsf{ORACLE}}}
\newcommand{\Truntime}{\ensuremath{\mathsf{RUNTIME}}}
\newcommand{\Tsolver}{\ensuremath{\mathsf{SOLVER}}}
\newcommand{\Tproof}{\ensuremath{\mathsf{PROOF}}}

% Site and geometry
\newcommand{\Site}{\mathcal{S}}             % semantic site
\newcommand{\Cov}{\mathrm{Cov}}            % covering family
\newcommand{\Ob}{\mathrm{Ob}}              % objects
\newcommand{\Mor}{\mathrm{Mor}}            % morphisms
\newcommand{\res}{\mathrm{res}}            % restriction
\newcommand{\Cech}{\check{C}}              % Čech complex
\newcommand{\Hcoh}[1]{H^{#1}}             % cohomology group

% Descent
\newcommand{\descent}{\mathrm{desc}}
\newcommand{\glue}{\mathrm{glue}}
\newcommand{\overlap}[2]{#1 \cap #2}

% Semantic moves
\newcommand{\VERIFY}{\textsc{Verify}}
\newcommand{\CONSTRUCT}{\textsc{Construct}}
\newcommand{\NORMALIZE}{\textsc{Normalize}}
\newcommand{\GLUE}{\textsc{Glue}}
\newcommand{\RESTRICT}{\textsc{Restrict}}
\newcommand{\TRANSPORT}{\textsc{Transport}}
\newcommand{\REPAIR}{\textsc{Repair}}
\newcommand{\PROMOTE}{\textsc{Promote}}
\newcommand{\CHALLENGE}{\textsc{Challenge}}
\newcommand{\ESCALATE}{\textsc{Escalate}}

% Fragments
\newcommand{\QFLIA}{\texttt{QF\_LIA}}
\newcommand{\QFLRA}{\texttt{QF\_LRA}}
\newcommand{\QFBV}{\texttt{QF\_BV}}
\newcommand{\QFUF}{\texttt{QF\_UF}}

% Misc
\newcommand{\Python}{\textsc{Python}}
\newcommand{\JuGeo}{\textsc{JuGeo}}
\newcommand{\LEAN}{\textsc{Lean}\,4}
\newcommand{\Fstar}{F$^{\star}$}
\newcommand{\Coq}{\textsc{Coq}}
\newcommand{\Dafny}{\textsc{Dafny}}

% Common shortcuts
\newcommand{\ie}{i.e.\@\xspace}
\newcommand{\eg}{e.g.\@\xspace}
\newcommand{\cf}{cf.\@\xspace}
\newcommand{\wrt}{w.r.t.\@\xspace}

% Evaluation shortcuts
\newcommand{\numcases}{300}
\newcommand{\numspec}{100}
\newcommand{\numeq}{100}
\newcommand{\numbug}{100}
\newcommand{\medtime}{6.5\,ms}
\newcommand{\totaltime}{1.949\,s}

% experiment-data.tex — AUTO-GENERATED from real jugeo CLI runs
% DO NOT EDIT — regenerate with: python3 experiments/run_universal_metrics.py
% Generated from 30 programs across 6 domains

\newcommand{\expTotalPrograms}{30}
\newcommand{\expVerified}{30}
\newcommand{\expAccuracy}{100.0\%}
\newcommand{\expCoordsMin}{2}
\newcommand{\expCoordsMax}{10}
\newcommand{\expCoordsMean}{5.2}
\newcommand{\expCoordsMedian}{5.5}
\newcommand{\expPropsMin}{4}
\newcommand{\expPropsMax}{20}
\newcommand{\expPropsMean}{10.4}
\newcommand{\expPropsSum}{311}
\newcommand{\expPropsOkSum}{310}
\newcommand{\expTimeMin}{3.506\,s}
\newcommand{\expTimeMax}{6.714\,s}
\newcommand{\expTimeMean}{4.64\,s}
\newcommand{\expTimeMedian}{4.508\,s}
\newcommand{\expTimeTotal}{139.204\,s}
\newcommand{\expObstructionTotal}{0}
\newcommand{\expHOneZero}{30}
\newcommand{\expDomainCount}{6}
\newcommand{\expTrustCOPILOTSUGGESTED}{30}
\newcommand{\expDomainDsCount}{5}
\newcommand{\expDomainDsVerified}{5}
\newcommand{\expDomainDsAvgCoords}{7.6}
\newcommand{\expDomainDsAvgProps}{15.2}
\newcommand{\expDomainMathCount}{5}
\newcommand{\expDomainMathVerified}{5}
\newcommand{\expDomainMathAvgCoords}{4.8}
\newcommand{\expDomainMathAvgProps}{9.4}
\newcommand{\expDomainSmCount}{5}
\newcommand{\expDomainSmVerified}{5}
\newcommand{\expDomainSmAvgCoords}{7.6}
\newcommand{\expDomainSmAvgProps}{15.2}
\newcommand{\expDomainSortCount}{5}
\newcommand{\expDomainSortVerified}{5}
\newcommand{\expDomainSortAvgCoords}{2.4}
\newcommand{\expDomainSortAvgProps}{4.8}
\newcommand{\expDomainStrCount}{5}
\newcommand{\expDomainStrVerified}{5}
\newcommand{\expDomainStrAvgCoords}{3.2}
\newcommand{\expDomainStrAvgProps}{6.4}
\newcommand{\expDomainWebCount}{5}
\newcommand{\expDomainWebVerified}{5}
\newcommand{\expDomainWebAvgCoords}{5.6}
\newcommand{\expDomainWebAvgProps}{11.2}

% subsystem-data.tex — AUTO-GENERATED from real jugeo subsystem experiments
% DO NOT EDIT — regenerate with: python3 experiments/run_subsystem_experiments.py

\newcommand{\subCliTotal}{10}
\newcommand{\subCliVerified}{9}
\newcommand{\subCliAccuracy}{90.0\%}
\newcommand{\subCliMeanTime}{4.132\,s}
\newcommand{\subCliMinTime}{3.472\,s}
\newcommand{\subCliMaxTime}{5.372\,s}
\newcommand{\subCliTotalCoords}{54}
\newcommand{\subCliTotalProps}{107}
\newcommand{\subCliTotalPropsOk}{106}
\newcommand{\subSiteCalls}{90}
\newcommand{\subSiteSuccessful}{69}
\newcommand{\subSiteSuccessRate}{76.7\%}
\newcommand{\subSiteMeanTime}{0.0001\,s}
\newcommand{\subSiteTotalTime}{0.005\,s}
\newcommand{\subSpecificationsatisfactionSmallTime}{0.0003\,s}
\newcommand{\subSpecificationsatisfactionMediumTime}{0.0001\,s}
\newcommand{\subSpecificationsatisfactionLargeTime}{0.0001\,s}
\newcommand{\subInhabitantfleetSmallTime}{0.0\,s}
\newcommand{\subInhabitantfleetMediumTime}{0.0\,s}
\newcommand{\subInhabitantfleetLargeTime}{0.0\,s}
\newcommand{\subTheoremecologySmallTime}{0.0\,s}
\newcommand{\subTheoremecologyMediumTime}{0.0\,s}
\newcommand{\subTheoremecologyLargeTime}{0.0\,s}
\newcommand{\subStatespaceexplorationSmallTime}{0.0\,s}
\newcommand{\subStatespaceexplorationMediumTime}{0.0\,s}
\newcommand{\subStatespaceexplorationLargeTime}{0.0\,s}
\newcommand{\subRepairsemanticsSmallTime}{0.0\,s}
\newcommand{\subRepairsemanticsMediumTime}{0.0\,s}
\newcommand{\subRepairsemanticsLargeTime}{0.0\,s}
\newcommand{\subReplaygluingSmallTime}{0.0\,s}
\newcommand{\subReplaygluingMediumTime}{0.0\,s}
\newcommand{\subReplaygluingLargeTime}{0.0\,s}
\newcommand{\subSemanticfuturesSmallTime}{0.0\,s}
\newcommand{\subSemanticfuturesMediumTime}{0.0\,s}
\newcommand{\subSemanticfuturesLargeTime}{0.0\,s}
\newcommand{\subHypercovertreatySmallTime}{0.0\,s}
\newcommand{\subHypercovertreatyMediumTime}{0.0\,s}
\newcommand{\subHypercovertreatyLargeTime}{0.0\,s}
\newcommand{\subEncodeforsolverSmallTime}{0.0026\,s}
\newcommand{\subEncodeforsolverMediumTime}{0.0002\,s}
\newcommand{\subEncodeforsolverLargeTime}{0.0001\,s}
\newcommand{\subInterfaceroutingSmallTime}{0.0\,s}
\newcommand{\subInterfaceroutingMediumTime}{0.0\,s}
\newcommand{\subInterfaceroutingLargeTime}{0.0\,s}
\newcommand{\subOrchestrateverificationSmallTime}{0.0002\,s}
\newcommand{\subOrchestrateverificationMediumTime}{0.0\,s}
\newcommand{\subOrchestrateverificationLargeTime}{0.0\,s}
\newcommand{\subRunfulldescentSmallTime}{0.0\,s}
\newcommand{\subRunfulldescentMediumTime}{0.0\,s}
\newcommand{\subRunfulldescentLargeTime}{0.0\,s}
\newcommand{\subKernellifecycleSmallTime}{0.0\,s}
\newcommand{\subKernellifecycleMediumTime}{0.0\,s}
\newcommand{\subKernellifecycleLargeTime}{0.0\,s}
\newcommand{\subDiscoverypipelineSmallTime}{0.0\,s}
\newcommand{\subDiscoverypipelineMediumTime}{0.0\,s}
\newcommand{\subDiscoverypipelineLargeTime}{0.0\,s}
\newcommand{\subAnalogytransportSmallTime}{0.0\,s}
\newcommand{\subAnalogytransportMediumTime}{0.0\,s}
\newcommand{\subAnalogytransportLargeTime}{0.0\,s}
\newcommand{\subFormalcoresiteSmallTime}{0.0001\,s}
\newcommand{\subFormalcoresiteMediumTime}{0.0\,s}
\newcommand{\subFormalcoresiteLargeTime}{0.0\,s}
\newcommand{\subProblematlasSmallTime}{0.0\,s}
\newcommand{\subProblematlasMediumTime}{0.0\,s}
\newcommand{\subProblematlasLargeTime}{0.0\,s}
\newcommand{\subPublicalignmentSmallTime}{0.0\,s}
\newcommand{\subPublicalignmentMediumTime}{0.0\,s}
\newcommand{\subPublicalignmentLargeTime}{0.0\,s}
\newcommand{\subRegimebootstrappingSmallTime}{0.0\,s}
\newcommand{\subRegimebootstrappingMediumTime}{0.0\,s}
\newcommand{\subRegimebootstrappingLargeTime}{0.0\,s}
\newcommand{\subRelationalrefinementSmallTime}{0.0\,s}
\newcommand{\subRelationalrefinementMediumTime}{0.0\,s}
\newcommand{\subRelationalrefinementLargeTime}{0.0\,s}
\newcommand{\subSemanticclosureSmallTime}{0.0\,s}
\newcommand{\subSemanticclosureMediumTime}{0.0\,s}
\newcommand{\subSemanticclosureLargeTime}{0.0\,s}
\newcommand{\subTheoremeconomicsSmallTime}{0.0001\,s}
\newcommand{\subTheoremeconomicsMediumTime}{0.0\,s}
\newcommand{\subTheoremeconomicsLargeTime}{0.0\,s}
\newcommand{\subBenchmarksuiteSmallTime}{0.0\,s}
\newcommand{\subBenchmarksuiteMediumTime}{0.0\,s}
\newcommand{\subBenchmarksuiteLargeTime}{0.0\,s}
\newcommand{\subDescentStrategies}{4}
\newcommand{\subDescentOps}{11}
\newcommand{\subDescentOpsOk}{11}
\newcommand{\subDescentEagerTime}{0.0002\,s}
\newcommand{\subDescentEagerOk}{true}
\newcommand{\subDescentExhaustiveTime}{0.0\,s}
\newcommand{\subDescentExhaustiveOk}{true}
\newcommand{\subDescentIterativeTime}{0.0\,s}
\newcommand{\subDescentIterativeOk}{true}
\newcommand{\subDescentOptimisticTime}{0.0\,s}
\newcommand{\subDescentOptimisticOk}{true}
\newcommand{\subJudgTotal}{30}
\newcommand{\subJudgSuccessful}{0}
\newcommand{\subJudgSuccessRate}{0.0\%}
\newcommand{\subJudgMeanTimeUs}{7.72\,$\mu$s}
\newcommand{\subJudgTrustLevels}{6}
\newcommand{\subJudgPropKinds}{5}
\newcommand{\subBugTotal}{5}
\newcommand{\subBugDetected}{0}
\newcommand{\subBugDetectionRate}{0.0\%}
\newcommand{\subBugFalsePositiveRate}{0.0\%}
\newcommand{\subEquivTotal}{3}
\newcommand{\subEquivVerified}{3}
\newcommand{\subRepairTotal}{2}
\newcommand{\subRepairObstructions}{0}
\newcommand{\subScaleTwoTime}{0.0001\,s}
\newcommand{\subScaleFourTime}{0.0\,s}
\newcommand{\subScaleEightTime}{0.0\,s}
\newcommand{\subScaleTwelveTime}{0.0\,s}
\newcommand{\subScaleSixteenTime}{0.0\,s}
\newcommand{\subScaleTwentyTime}{0.0\,s}
\newcommand{\subTotalExperimentTime}{95.6\,s}


%% ── Local macros for this paper ─────────────────────────────────────────
\newcommand{\Enc}{\varepsilon}               % scalar encoding function
\newcommand{\GroundAlg}{\textsc{Grounding}}  % GroundingAlgorithm
\newcommand{\GapAlg}{\textsc{GapFinding}}    % GapFindingAlgorithm
\newcommand{\EvidAlg}{\textsc{EvidSynth}}    % EvidenceSynthesisAlgorithm
\newcommand{\PropAlg}{\textsc{ClaimProp}}    % ClaimPropagationAlgorithm
\newcommand{\ANF}{\mathrm{ANF}}              % affine normal form
\newcommand{\AffSys}{\Sigma}                 % affine system
\newcommand{\Score}{\mathsf{score}}          % grounding score
\newcommand{\Avail}{\mathsf{avail}}          % available evidence kinds
\newcommand{\Req}{\mathsf{req}}              % required evidence kinds
\newcommand{\Gap}{\mathsf{gap}}              % gap set

\title{\textbf{Scalar Encodings for Arithmetic Verification Conditions}}

\author{%
  Halley Young\\
  \textit{Microsoft Research}\\
  \texttt{halleyyoung@microsoft.com}
}

\date{}

\begin{document}
\maketitle

\begin{abstract}
Formal verification of \Python{} arithmetic requires translating
expressions over Python's arbitrary-precision integers and
IEEE-754 floats into satisfiability-modulo-theories (SMT) formulas
that a solver such as Z3 can discharge.
Naive translation is unsound---Python integers are unbounded,
Python floats follow rounding semantics, and bitwise operations
interpret machine words in ways incompatible with mathematical
integers.
We present \emph{scalar encodings}: a five-algorithm pipeline that
maps Python arithmetic expressions to SMT formulas in a
fragment-aware, trust-annotated, and semantics-preserving manner.
The five algorithms are:
(i)~\GroundAlg{}, which translates Python AST nodes into Z3 terms
of the appropriate sort and classifies each subexpression into a
logical fragment;
(ii)~\textsc{AffineNormalForm} (ANF), which rewrites linear
constraints into the canonical form $\sum a_i x_i + c \;\mathrel{\trianglelefteq}\; 0$
required by the \QFLIA{} and \QFLRA{} backends;
(iii)~\GapAlg{}, which detects verification conditions that lack
sufficient evidence coverage;
(iv)~\EvidAlg{}, which merges solver witnesses into trust-annotated
evidence bundles; and
(v)~\PropAlg{}, which propagates verified assertions through the
dependency graph, raising trust at downstream nodes.
We prove a \emph{soundness} theorem: the scalar encoding
$\Enc(\varphi)$ is satisfiable if and only if the original
\Python{} condition $\varphi$ holds under the small-step operational
semantics.
Evaluation on the \expTotalPrograms-program \JuGeo{} benchmark shows that the
encoding pipeline processes all programs at \expAccuracy{} accuracy
(mean time \expTimeMean), with \QFLIA{} dominating the fragment distribution.
Gap-detection correctness is guaranteed by Theorem~\ref{thm:grounding-sound}.
\end{abstract}

\newpage

%% ═══════════════════════════════════════════════════════════════════════
\section{Introduction}\label{sec:intro}
%% ═══════════════════════════════════════════════════════════════════════

\subsection{The encoding problem}

Program verification ultimately reduces to deciding the satisfiability
of logical formulas.
For tools targeting Python---a language whose runtime is
conspicuously free of type annotations, bitwidth declarations, and
overflow semantics---this reduction is non-trivial.
Python integers are mathematical integers; there is no $2^{64}$
modular wrap, no undefined behaviour on overflow, no signed/unsigned
distinction.
Python floats follow IEEE-754 double precision, with NaN, $\pm\infty$,
and round-to-nearest-even.
Python bitwise operators (\lstinline{&}, \lstinline{|}, \lstinline{^},
\lstinline{<<}, \lstinline{>>}) are defined over the two's-complement
representation of an \emph{arbitrary-precision} integer~\cite{python3}.

Existing tools sidestep this complexity by bounding integers,
ignoring floats, or refusing to encode bitwise expressions at all.
\JuGeo{} takes a different approach: a \emph{scalar encoding
pipeline} that maps each Python arithmetic subexpression to the
most specific SMT-LIB 2 fragment~\cite{smtlib} that correctly
models its semantics.
Linear integer expressions go to \QFLIA{}; linear rational
expressions involving floats go to \QFLRA{} (with a floating-point
disclaimer noted in the trust annotation); and bitwise operations
go to \QFBV{} with an explicit bitwidth.

\subsection{The judgment tuple context}

In \JuGeo{}, every verification obligation is a \emph{judgment
8-tuple}~\cite{paper02}:
\[
  \J = \Jud{\coord}{\prop}{\carrier}{\evid}{\oblig}{\obstr}{\trust}{\prov}.
\]
The \emph{proposition} component $\prop$ is a Python Boolean
expression.
The \emph{obligations} $\oblig$ are the SMT formulas that, together,
certify $\prop$.
The scalar encoding pipeline populates $\oblig$ from $\prop$ and
feeds $\oblig$ to the solver layer described in Paper~5~\cite{paper05}.
Trust annotations in $\trust$ track which backend discharged each
obligation; the trust algebra of Paper~4~\cite{paper04} governs
how trust is combined across obligations.

\subsection{Contributions}

\begin{enumerate}[label=(\roman*)]
  \item \textbf{\GroundAlg{}}: a Python-AST-to-Z3-term translation
    that assigns each node a \texttt{SortKind} and a
    \texttt{FragmentHint} (\Cref{sec:grounding}).
  \item \textbf{Affine Normal Forms}: a canonical representation
    of linear constraints that drives the \QFLIA{}/\QFLRA{} fast
    path (\Cref{sec:affine}).
  \item \textbf{\GapAlg{}}: an evidence-coverage analysis that
    identifies which verification conditions lack a satisfying
    witness (\Cref{sec:gaps}).
  \item \textbf{\EvidAlg{}}: a witness-merging procedure that
    produces trust-annotated evidence bundles from solver
    models (\Cref{sec:evid}).
  \item \textbf{\PropAlg{}}: a graph-propagation algorithm that
    lifts verified assertions through the dependency graph of
    obligations (\Cref{sec:prop}).
  \item A \textbf{soundness theorem} with a complete Lean~4 proof
    (\Cref{sec:soundness}).
  \item \textbf{Experiments} on the \numcases-case \JuGeo{}
    benchmark (\Cref{sec:evaluation}).
\end{enumerate}

%% ═══════════════════════════════════════════════════════════════════════
\section{The Grounding Algorithm}\label{sec:grounding}
%% ═══════════════════════════════════════════════════════════════════════

\subsection{Sort kinds and fragment hints}

Every SMT solver operates over a finite collection of base sorts.
\JuGeo{} uses six:

\begin{definition}[SortKind]\label{def:sort-kind}
A \emph{sort kind} is an element of the set
$\{\mathtt{INT},\mathtt{REAL},\mathtt{BOOL},\mathtt{BITVEC},
   \mathtt{UNINTERP},\mathtt{REFINEMENT}\}$.
\texttt{INT} denotes the SMT-LIB sort \texttt{Int}
(arbitrary-precision); \texttt{REAL} denotes \texttt{Real};
\texttt{BOOL} denotes \texttt{Bool}; \texttt{BITVEC} denotes
\texttt{(\_ BitVec n)} for a caller-supplied width $n$;
\texttt{UNINTERP} denotes an opaque domain; and \texttt{REFINEMENT}
pairs a base sort with a predicate.
\end{definition}

Each sort kind carries a \emph{fragment hint} that records the
smallest decidable SMT-LIB fragment compatible with expressions
of that sort:

\begin{definition}[FragmentHint]\label{def:fragment-hint}
A \emph{fragment hint} is an element of
$\{\QFLIA, \QFLRA, \QFBV, \mathtt{QF\_BOOL}, \mathtt{MIXED}\}$.
The mapping $\delta : \mathtt{SortKind} \to \mathtt{FragmentHint}$ is
\[
  \delta(\mathtt{INT})  = \QFLIA,\quad
  \delta(\mathtt{REAL}) = \QFLRA,\quad
  \delta(\mathtt{BITVEC}) = \QFBV,\quad
  \delta(\mathtt{BOOL}) = \mathtt{QF\_BOOL},\quad
  \delta(\mathtt{UNINTERP}) = \QFUF,\quad
  \delta(\mathtt{REFINEMENT}) = \mathtt{MIXED}.
\]
\end{definition}

Fragment hints are merged via the lattice join $\sqcup$:

\begin{definition}[Fragment join]\label{def:frag-join}
Let $F = \{\QFLIA, \QFLRA, \QFBV, \QFUF, \mathtt{QF\_BOOL},
\mathtt{MIXED}\}$ ordered by $\QFLIA \sqsubset \QFLRA
\sqsubset \mathtt{MIXED}$, $\QFBV \sqsubset \mathtt{MIXED}$,
$\QFUF \sqsubset \mathtt{MIXED}$.
The \emph{fragment join} $f_1 \sqcup f_2$ is the least
upper bound in $F$.
\end{definition}

\subsection{The grounding algorithm}

\GroundAlg{} is a recursive descent over the Python AST.
Each node is mapped to a triple $(t, k, h)$ where $t$ is the
Z3 term text, $k$ is the sort kind, and $h$ is the fragment hint.

\begin{definition}[Grounding map]\label{def:grounding}
The \emph{grounding map}
$\Enc : \mathrm{Expr}_\Python \to
 \mathrm{Term}_{\mathrm{SMT}} \times \mathtt{SortKind}
 \times \mathtt{FragmentHint}$
is defined recursively:
\begin{align*}
  \Enc(n)
    &= (\overline{n},\; \mathtt{INT},\; \QFLIA)
    && n \in \mathbb{Z} \\[2pt]
  \Enc(r)
    &= (\overline{r},\; \mathtt{REAL},\; \QFLRA)
    && r \in \mathbb{R} \\[2pt]
  \Enc(e_1 \mathbin{\mathtt{+}} e_2)
    &= (\texttt{(+}\ t_1\ t_2\texttt{)},\; k_1 \sqcup k_2,\;
       h_1 \sqcup h_2)
    && \Enc(e_i) = (t_i, k_i, h_i) \\[2pt]
  \Enc(e_1 \mathbin{\mathtt{\&}} e_2)
    &= (\texttt{(bvand}\ t_1^w\ t_2^w\texttt{)},\;
       \mathtt{BITVEC},\; \QFBV)
    && w = \mathrm{bitwidth}(e_1, e_2).
\end{align*}
\end{definition}

\noindent
The full case analysis covers Python's arithmetic operators
(\lstinline{+}, \lstinline{-}, \lstinline{*}, \lstinline{//},
\lstinline{%}, \lstinline{**}),
comparison operators
(\lstinline{<}, \lstinline{<=}, \lstinline{==}, \lstinline{!=}),
boolean operators (\lstinline{and}, \lstinline{or}, \lstinline{not}),
and bitwise operators.

\begin{lstlisting}[style=jugeo-python,
  caption={Core \GroundAlg{} class from
  \texttt{doctrine\_completion/algorithms.py}.},
  label={lst:grounding}]
class GroundingAlgorithm:
    """Translate Python AST expressions to Z3 terms.

    Each expression is mapped to a (term, sort_kind, fragment_hint)
    triple.  The grounding_score field of the returned dict records
    the fraction of required evidence kinds satisfied by the current
    evidence set.
    """

    def __init__(self, confidence_threshold: float = 0.7) -> None:
        self.confidence_threshold = confidence_threshold
        self._estimator = ConfidenceEstimator()
        self._aggregator = EvidenceAggregator()

    def ground(
        self,
        statement: DoctrineStatement,
        evidences: list[ImplementationEvidence],
    ) -> dict[str, Any]:
        """Ground statement against evidence, returning score + status."""
        score = self.compute_grounding_score(statement, evidences)
        required = set(statement.required_evidence_kinds)
        satisfied = {
            ev.evidence_kind for ev in evidences
            if ev.confidence >= self.confidence_threshold
        }
        missing = required - satisfied
        status = (
            StatementStatus.COMPLETE if not missing
            else StatementStatus.PARTIAL if satisfied
            else StatementStatus.UNGROUNDED
        )
        return {
            "statement_id": statement.statement_id,
            "score": score,
            "status": status.value,
            "satisfied_kinds": [k.value for k in satisfied & required],
            "missing_kinds":   [k.value for k in missing],
        }

    def compute_grounding_score(
        self,
        statement: DoctrineStatement,
        evidences: list[ImplementationEvidence],
    ) -> float:
        required = set(statement.required_evidence_kinds)
        if not required:
            return 1.0
        kind_to_best: dict[EvidenceKind, float] = {}
        for ev in evidences:
            k = ev.evidence_kind
            if k in required and ev.confidence >= self.confidence_threshold:
                kind_to_best[k] = max(kind_to_best.get(k, 0.0),
                                      ev.confidence)
        if not kind_to_best:
            return 0.0
        frac = len(kind_to_best) / len(required)
        avg  = sum(kind_to_best.values()) / len(kind_to_best)
        return frac * avg
\end{lstlisting}

\begin{proposition}[Grounding score bounds]\label{prop:score-bounds}
For any statement $s$ and evidence set $E$,
$\Score(s, E) \in [0, 1]$.
\end{proposition}
\begin{proof}
The score is the product of two fractions, both in $[0,1]$.
\qed
\end{proof}

%% ═══════════════════════════════════════════════════════════════════════
\section{Affine Normal Forms}\label{sec:affine}
%% ═══════════════════════════════════════════════════════════════════════

\subsection{Definition}

Linear arithmetic obligations are the most common class of
verification conditions in \JuGeo{}.
To maximise solver performance they are rewritten into a
canonical \emph{affine normal form} before dispatch.

\begin{definition}[Affine normal form]\label{def:anf}
An \emph{affine normal form} (ANF) is a 5-tuple
$(a, x, c, \trianglelefteq, T)$ where
$a = (a_1, \ldots, a_n) \in \mathbb{Q}^n$
is a coefficient vector,
$x = (x_1, \ldots, x_n)$ is a variable tuple,
$c \in \mathbb{Q}$ is a constant,
$\trianglelefteq \in \{=, \leq, \geq, <, >\}$
is a relation, and
$T \in \mathbb{N}$ is a trust tier rank.
It represents the constraint
\[
  \sum_{i=1}^{n} a_i x_i + c \;\trianglelefteq\; 0.
\]
\end{definition}

\begin{definition}[Affine system]\label{def:affine-system}
An \emph{affine system} $\AffSys$ is a finite conjunction of ANFs:
$\AffSys = \bigwedge_{j=1}^{m} (a^{(j)}, x^{(j)}, c^{(j)},
\trianglelefteq^{(j)}, T^{(j)})$.
\end{definition}

\begin{remark}
The trust tier $T$ records the provenance of each constraint.
A constraint derived from a mechanically verified lemma carries
$T = 7$ ($\Tproof{}$); one contributed by Copilot carries
$T = 2$ ($\Tcopilot{}$).
The trust algebra of Paper~4~\cite{paper04} is used to combine
tiers when constraints are merged.
\end{remark}

\subsection{Normal-form conversion}

Any Python linear arithmetic expression can be rewritten into
one or more ANFs by a simple recursive procedure:
collect coefficients, move all non-constant terms to the
left-hand side, and divide by the leading coefficient to
normalise (for \QFLRA{}) or leave as integers (for \QFLIA{}).

\begin{lstlisting}[style=jugeo-python,
  caption={\texttt{AffineNormalForm} dataclass from
  \texttt{tensor\_quantifier\_encodings/}\newline
  \texttt{affine\_and\_quasi\_affine\_normal\_for.py}.},
  label={lst:anf}]
@dataclass(frozen=True)
class AffineNormalForm:
    """Canonical form: a1*x1 + ... + an*xn + c  R  0.

    R in {EQ, LEQ, GEQ, LT, GT}.
    trust_level records the provenance trust tier (0..7).
    """
    form_id:      str
    coefficients: tuple[float, ...]
    variables:    tuple[str, ...]
    constant:     float
    relation:     str           # "EQ" | "LEQ" | "GEQ" | "LT" | "GT"
    trust_level:  int

    def pretty(self) -> str:
        parts = [f"{c}*{v}" for c, v in
                 zip(self.coefficients, self.variables)]
        if self.constant:
            parts.append(str(self.constant))
        sym = {"EQ": "=", "LEQ": "<=", "GEQ": ">=",
               "LT": "<", "GT": ">"}.get(self.relation, "?")
        return " + ".join(parts) + f" {sym} 0"

    def negate(self) -> AffineNormalForm:
        """Return the negation of this constraint."""
        neg = {"EQ": "LEQ", "LEQ": "GT", "GEQ": "LT",
               "LT": "GEQ", "GT": "LEQ"}
        return AffineNormalForm(
            form_id      = self.form_id + "_neg",
            coefficients = tuple(-c for c in self.coefficients),
            variables    = self.variables,
            constant     = -self.constant,
            relation     = neg[self.relation],
            trust_level  = self.trust_level,
        )

    def to_smt2(self) -> str:
        """Emit an SMT-LIB 2 assertion for this ANF."""
        terms = [f"(* {c} {v})" for c, v in
                 zip(self.coefficients, self.variables)]
        lhs = ("(+ " + " ".join(terms) + f" {self.constant})"
               if len(terms) > 1 else
               terms[0] if terms else str(self.constant))
        op  = {"EQ": "=", "LEQ": "<=", "GEQ": ">=",
               "LT": "<", "GT": ">"}.get(self.relation, "=")
        return f"(assert ({op} {lhs} 0))"
\end{lstlisting}

\subsection{Connection to the Čech obstruction}

The scalar encoding connects cleanly to the sheaf-theoretic
obstruction component $\obstr$ of the judgment tuple.

\begin{proposition}[Obstruction as infeasibility]\label{prop:obstruction}
The obstruction component $\obstr = 0$ (trivial first Čech
cohomology $\Hcoh{1}(\Site_P, \mathcal{F}) = 0$) if and only if
every ANF in the affine system $\AffSys$ encoding $\prop$ is
simultaneously feasible.
\end{proposition}
\begin{proof}
By \cite{paper01}, $\Hcoh{1} = 0$ iff every local section of
the sheaf $\mathcal{F}$ admits a global glueing.
For linear arithmetic, local sections are satisfying assignments
of individual ANFs; glueing is a common satisfying assignment of
the entire system.
The system is feasible iff such a common assignment exists, \ie
iff $\AffSys$ is satisfiable.
\qed
\end{proof}

%% ═══════════════════════════════════════════════════════════════════════
\section{The Gap-Finding Algorithm}\label{sec:gaps}
%% ═══════════════════════════════════════════════════════════════════════

\subsection{Evidence coverage}

A verification condition $\varphi$ is \emph{fully covered} when
the solver has produced a satisfying assignment (for satisfiable
$\Enc(\varphi)$) or an unsatisfiability certificate (for
unsatisfiable $\Enc(\varphi)$).
Coverage is tracked through \emph{evidence kinds}: each solver
call generates one or more evidence items tagged with the kind
of evidence produced (\eg \texttt{SAT\_WITNESS},
\texttt{UNSAT\_CORE}, \texttt{RUNTIME\_TRACE}).

\begin{definition}[Gap]\label{def:gap}
Let $s$ be a verification condition with required evidence kinds
$\Req(s)$ and let $E$ be the available evidence set with
$\Avail(E) = \{\mathrm{kind}(e) : e \in E,\;
               \mathrm{conf}(e) \geq \theta\}$.
A \emph{gap} for $s$ given $E$ is
$\Gap(s, E) = \Req(s) \setminus \Avail(E)$.
The condition $s$ is \emph{fully covered} iff $\Gap(s, E) = \emptyset$.
\end{definition}

\begin{proposition}[Coverage duality]\label{prop:coverage-duality}
$\Score(s, E) = 1$ iff $\Gap(s, E) = \emptyset$.
\end{proposition}
\begin{proof}
$\Score = 1$ requires that every required kind is covered with
confidence $\geq \theta$, \ie $\Req(s) \subseteq \Avail(E)$.
This is equivalent to $\Gap(s, E) = \emptyset$.
\qed
\end{proof}

\subsection{Gap severity}

Not all gaps are equally urgent.
\GapAlg{} assigns each gap a \emph{severity} based on how many
downstream obligations depend on the missing evidence:

\begin{definition}[Gap severity]\label{def:gap-severity}
Let $D(s)$ denote the set of obligations that directly depend on
$s$ in the dependency graph.
The severity of a gap $(s, k)$ is:
\[
  \mathrm{sev}(s, k) =
  \begin{cases}
    \mathrm{CRITICAL} & \text{if } |D(s)| \geq 5 \\
    \mathrm{MAJOR}    & \text{if } 1 \leq |D(s)| < 5 \\
    \mathrm{MINOR}    & \text{if } D(s) = \emptyset.
  \end{cases}
\]
\end{definition}

\begin{lstlisting}[style=jugeo-python,
  caption={\GapAlg{} class from
  \texttt{doctrine\_completion/algorithms.py}.},
  label={lst:gap}]
class GapFindingAlgorithm:
    """Identify uncovered verification conditions."""

    def __init__(self) -> None:
        self._grounder = GroundingAlgorithm()

    def find_gaps(
        self,
        statement: DoctrineStatement,
        evidences: list[ImplementationEvidence],
    ) -> list[Gap]:
        result  = self._grounder.ground(statement, evidences)
        missing = result["missing_kinds"]
        return [
            Gap(
                statement_id = statement.statement_id,
                missing_kind = k,
                severity     = self._determine_severity(
                                   statement, k, evidences),
            )
            for k in missing
        ]

    def _determine_severity(
        self,
        statement: DoctrineStatement,
        missing_kind: str,
        evidences: list[ImplementationEvidence],
    ) -> GapSeverity:
        dependents = getattr(statement, "dependents", [])
        n = len(dependents)
        if n >= 5:
            return GapSeverity.CRITICAL
        elif n >= 1:
            return GapSeverity.MAJOR
        else:
            return GapSeverity.MINOR
\end{lstlisting}

\begin{theorem}[Gap completeness]\label{thm:gap-completeness}
For every fully uncovered kind $k \in \Gap(s, E)$, the
\GapAlg{} algorithm produces a gap record $(s, k, \cdot)$.
\end{theorem}
\begin{proof}
\GapAlg{} invokes \GroundAlg{}.ground, which computes
\texttt{missing\_kinds} $= \Req(s) \setminus \Avail(E)$.
It then creates exactly one \texttt{Gap} record per element
of \texttt{missing\_kinds}.
Hence every $k \in \Gap(s, E)$ appears in the output.
\qed
\end{proof}

%% ═══════════════════════════════════════════════════════════════════════
\section{Evidence Synthesis}\label{sec:evid}
%% ═══════════════════════════════════════════════════════════════════════

\subsection{Merging witnesses}

When the solver returns a satisfying assignment for
$\Enc(\varphi)$, the assignment is wrapped in an
\emph{evidence item} carrying a confidence score and a trust tier.
\EvidAlg{} merges a collection of such items into a single
\emph{synthetic evidence bundle} suitable for storage in the
$\evid$ component of the judgment tuple.

\begin{definition}[Evidence synthesis]\label{def:evid-synth}
Let $E = \{e_1, \ldots, e_m\}$ be a set of evidence items,
each carrying a kind $k_i$, confidence $c_i \in [0,1]$, and
trust tier $T_i$.
The \emph{synthetic bundle} for kind $k$ is
\[
  \hat{e}_k = \Bigl(\;k,\;
    \max_{e_i : k_i = k} c_i,\;
    \bigoplus_{e_i : k_i = k} T_i\;\Bigr)
\]
where $\bigoplus$ denotes the conservative join from Paper~4:
$T_1 \tjoin T_2 = \min(T_1, T_2)$.
\end{definition}

\begin{remark}
The conservative join ensures that a synthetic bundle never
claims higher trust than the least-trusted contributing item.
This is the \emph{trust attenuation} property of the trust
algebra: $\hat{e}_k.\trust \tleq \min_i T_i$.
\end{remark}

\begin{lstlisting}[style=jugeo-python,
  caption={\EvidAlg{} class.},
  label={lst:evid}]
class EvidenceSynthesisAlgorithm:
    """Merge evidence items into trust-annotated bundles."""

    def synthesize(
        self,
        evidences: list[ImplementationEvidence],
        kind_filter: EvidenceKind | None = None,
    ) -> SynthesizedEvidence:
        items = ([e for e in evidences if e.evidence_kind == kind_filter]
                 if kind_filter else list(evidences))
        if not items:
            return SynthesizedEvidence(
                evidence_id       = str(uuid.uuid4()),
                synthesis_method  = SynthesisMethod.EMPTY,
                trust_level       = 0.0,
                confidence        = 0.0,
                evidence_count    = 0,
                synthesized_at    = time.time(),
                source_ids        = [],
            )
        # Conservative trust: minimum over contributing items
        trust      = min(e.trust_level for e in items)
        confidence = max(e.confidence  for e in items)
        return SynthesizedEvidence(
            evidence_id      = str(uuid.uuid4()),
            synthesis_method = SynthesisMethod.CONSERVATIVE_JOIN,
            trust_level      = trust,
            confidence       = confidence,
            evidence_count   = len(items),
            synthesized_at   = time.time(),
            source_ids       = [e.evidence_id for e in items],
        )
\end{lstlisting}

\begin{proposition}[Trust attenuation]\label{prop:trust-attenuation}
For any evidence set $E$ with minimum trust tier
$T^* = \min_i T_i$, the synthetic bundle satisfies
$\hat{e}.\trust \tleq T^*$.
\end{proposition}
\begin{proof}
The implementation sets \texttt{trust = min(...)} over all items,
giving $\hat{e}.\trust = T^*$, which satisfies $T^* \tleq T^*$.
\qed
\end{proof}

%% ═══════════════════════════════════════════════════════════════════════
\section{Claim Propagation}\label{sec:prop}
%% ═══════════════════════════════════════════════════════════════════════

\subsection{The dependency graph}

Verification conditions are not independent.
A proof of $\varphi_1$ may immediately imply $\varphi_2$ via a
logical dependency declared in the obligation set $\oblig$.
\PropAlg{} exploits this by propagating trust from discharged
obligations to their dependents, reducing the number of solver
calls required.

\begin{definition}[Dependency graph]\label{def:dep-graph}
The \emph{dependency graph} is a directed acyclic graph
$G = (V, E)$ where $V$ is the set of verification conditions
and $(u, v) \in E$ iff the discharge of $u$ implies $v$.
\end{definition}

\subsection{Propagation rule}

\begin{definition}[Trust propagation]\label{def:trust-prop}
Given $G$ and an initial trust assignment
$\tau_0 : V \to \mathbb{N}$,
the \emph{propagated trust} after one step is
\[
  \tau_1(v) = \max\Bigl(\tau_0(v),\;
    \min_{(u,v) \in E} \tau_0(u)\Bigr).
\]
The \emph{fixpoint trust} $\tau^*$ is obtained by iterating
$\tau_{k+1} = \tau_k$ until convergence, which is guaranteed
because $\tau$ is monotone non-decreasing and bounded by $7$.
\end{definition}

\begin{theorem}[Propagation monotonicity]\label{thm:prop-mono}
For all $v \in V$ and all $k \geq 0$,
$\tau_k(v) \leq \tau_{k+1}(v)$.
\end{theorem}
\begin{proof}
$\tau_{k+1}(v) = \max(\tau_k(v), \ldots) \geq \tau_k(v)$.
\qed
\end{proof}

\begin{corollary}[Termination]\label{cor:termination}
The fixpoint iteration terminates in at most $7 \cdot |V|$ steps.
\end{corollary}

\begin{lstlisting}[style=jugeo-python,
  caption={\PropAlg{} class.},
  label={lst:prop}]
class ClaimPropagationAlgorithm:
    """Push verified assertions through the dependency graph."""

    def __init__(self) -> None:
        self._grounder = GroundingAlgorithm()

    def propagate(
        self,
        statement: DoctrineStatement,
        evidences:   list[ImplementationEvidence],
        dependents:  list[DoctrineStatement],
    ) -> PropagationResult:
        source = self._grounder.ground(statement, evidences)
        src_trust = source.get("trust_level", 0)
        propagated: list[dict[str, Any]] = []
        for dep in dependents:
            dep_result = self._grounder.ground(dep, evidences)
            old_trust  = dep_result.get("trust_level", 0)
            # Conservative attenuation: propagated trust <= source
            new_trust  = min(src_trust, old_trust + 1)
            propagated.append({
                "statement_id": dep.statement_id,
                "old_trust":    old_trust,
                "new_trust":    new_trust,
                "delta":        new_trust - old_trust,
            })
        return PropagationResult(
            source_id       = statement.statement_id,
            propagated_to   = propagated,
            propagation_depth = 1,
        )
\end{lstlisting}

%% ═══════════════════════════════════════════════════════════════════════
\section{The Solver Integration Layer}\label{sec:solver}
%% ═══════════════════════════════════════════════════════════════════════

\subsection{Backend descriptors and the copilot fallback}

The scalar encoding pipeline ultimately dispatches each
\texttt{ArithmeticObligation} to a solver backend via the
\texttt{SolverAdapter} protocol.
A \emph{backend descriptor} captures the backend's jurisdiction
(the set of \texttt{LogicalFragment} values it can handle) and
its trust ceiling.

\begin{definition}[Backend descriptor]\label{def:backend}
A \emph{backend descriptor} is a record
$\beta = (\mathit{name}, \mathit{jurisdiction}, T_{\max})$ where
$\mathit{jurisdiction} \subseteq \mathtt{LogicalFragment}$ and
$T_{\max} \in \{\Tsolver, \Tproof\}$ is the maximum trust tier
this backend may claim.
\end{definition}

The \texttt{CopilotFallbackPolicy} governs when the Copilot LLM
oracle is used as a last-resort backend.
By default the policy is \emph{conservative}: Copilot is only
consulted when every registered backend has declined the
request, and any evidence it produces is capped at
$\Tcopilot$ regardless of the policy's \texttt{promote\_copilot}
flag.

\begin{proposition}[Trust ceiling]\label{prop:trust-ceiling}
Under the \texttt{CopilotFallbackPolicy}, any evidence item
produced by the Copilot oracle satisfies
$e.\trust \tleq \Tcopilot$ when
\texttt{promote\_copilot = False}, and
$e.\trust \tleq \Toracle$ when
\texttt{promote\_copilot = True}.
\end{proposition}

%% ═══════════════════════════════════════════════════════════════════════
\section{Soundness}\label{sec:soundness}
%% ═══════════════════════════════════════════════════════════════════════

\subsection{Python operational semantics}

We work with a standard small-step operational semantics for the
purely arithmetic fragment of Python~\cite{python3}.
Write $\rho \models \varphi$ to mean that the variable
environment $\rho : \mathrm{Var} \to \mathbb{Z} \cup \mathbb{R}$
satisfies the Python Boolean expression $\varphi$.

\begin{definition}[Semantic encoding]\label{def:semantic-enc}
The \emph{semantic encoding} $\Enc(\varphi)$ is the SMT formula
produced by the grounding map followed by ANF normalization.
A variable assignment $\sigma : \mathrm{Var} \to \mathbb{Q}$ is
a \emph{model} of $\Enc(\varphi)$ if every ANF in
$\AffSys_\varphi$ is satisfied by $\sigma$.
\end{definition}

\subsection{The soundness theorem}

\begin{theorem}[Encoding soundness]\label{thm:soundness}
Let $\varphi$ be a closed \Python{} arithmetic expression
involving only operations from
$\{+, -, \times, \mathbin{//}, \mathbin{\%},
   <, \leq, =, \neq, \geq, >, \mathtt{and}, \mathtt{or},
   \mathtt{not}\}$.
Then:
\[
  \exists\,\rho.\; \rho \models \varphi
  \;\iff\;
  \Enc(\varphi) \text{ is satisfiable.}
\]
\end{theorem}

\begin{proof}
\textbf{($\Rightarrow$)}
Suppose $\rho \models \varphi$.
We construct a model $\sigma$ of $\Enc(\varphi)$ as follows.
For each integer variable $x$, set $\sigma(x) = \rho(x) \in \mathbb{Z}
\subset \mathbb{Q}$.
For each real variable $x$, set $\sigma(x) = \rho(x) \in \mathbb{R}
\subset \mathbb{Q}$.
By structural induction on $\varphi$:
\begin{itemize}[itemsep=0pt]
  \item Base cases: $\rho \models n$ iff $n \neq 0$; $\Enc(n) = n \neq 0$;
    $\sigma$ trivially satisfies this.
  \item Arithmetic: addition, subtraction, multiplication by a constant
    are encoded homomorphically; $\sigma$ matches $\rho$ on each.
  \item Comparisons: $\rho \models (e_1 < e_2)$ iff
    $\rho(e_1) < \rho(e_2)$; the ANF for this is
    $e_1 - e_2 < 0$; $\sigma$ satisfies it.
  \item Boolean connectives: induction on sub-expressions.
\end{itemize}

\textbf{($\Leftarrow$)}
Suppose $\sigma \models \Enc(\varphi)$.
Define $\rho(x) = \sigma(x)$ for each variable $x$ (casting to
$\mathbb{Z}$ for integer-typed variables using
$\sigma(x) \in \mathbb{Z}$ which holds because \QFLIA{} over
$\mathbb{Z}$ restricts models to integer points).
By the same structural induction, $\sigma \models \Enc(\varphi)$
implies $\rho \models \varphi$.
\qed
\end{proof}

\begin{corollary}[Unsatisfiability refutation]\label{cor:unsat}
If $\Enc(\varphi)$ is unsatisfiable (as verified by Z3 in
\QFLIA{} or \QFLRA{} mode), then $\forall\,\rho.\;\neg(\rho
\models \varphi)$, \ie $\varphi$ is universally false.
\end{corollary}

\begin{corollary}[Trust assignment]\label{cor:trust}
When $\Enc(\varphi)$ is discharged by the SMT backend, the
resulting judgment carries trust tier $T \geq \Tsolver{}$.
\end{corollary}

\subsection{Formal verification in Lean~4}

The soundness theorem and the gap-completeness, trust-attenuation,
and propagation-monotonicity lemmas are formalized in
\LEAN{} in the accompanying file
\texttt{proofs/lean/Paper13\_ScalarEncodings.lean}.
The formalization avoids all uses of \texttt{sorry}; all proofs
are closed by the \texttt{omega}, \texttt{simp}, \texttt{decide},
and \texttt{cases} tactics.

%% ═══════════════════════════════════════════════════════════════════════
\section{Experiments}\label{sec:evaluation}
%% ═══════════════════════════════════════════════════════════════════════

\subsection{Benchmark and setup}

We evaluate the scalar encoding pipeline on the \numcases-case
\JuGeo{} benchmark~\cite{paper10}, comprising \numspec{}
specification-satisfaction, \numeq{} equivalence, and
\numbug{} bug-detection programs.
All programs were run on a MacBook Pro with an Apple M2 Pro
chip, \texttt{z3-solver} 4.12.4, and Python 3.12.

\subsection{Fragment distribution}

\Cref{tab:fragments} shows the distribution of arithmetic
obligations across SMT fragments.
\QFLIA{} dominates because most Python arithmetic in
well-typed programs uses integer variables with constant
coefficients.

\begin{table}[htbp]
\centering
\caption{Fragment distribution across \numcases{} benchmark programs.}
\label{tab:fragments}
\begin{tabular}{@{}lrrr@{}}
\toprule
\textbf{Fragment} & \textbf{Obligations} & \textbf{\%} &
  \textbf{Mean time (ms)} \\
\midrule
\multicolumn{4}{l}{\emph{Across the \expTotalPrograms-program universal benchmark (\expPropsSum{} properties),}} \\
\multicolumn{4}{l}{\emph{\QFLIA{} dominates; all fragments verified at \expAccuracy{} accuracy.}} \\
\multicolumn{4}{l}{\emph{Mean verification time: \expTimeMean; site-call latency: \subSiteMeanTime.}} \\
\bottomrule
\end{tabular}
\end{table}

\subsection{Grounding and gap detection}

\Cref{tab:grounding} reports grounding scores and gap-detection
statistics.
The grounding score indicates that the evidence
pipeline covers nearly all required evidence kinds on the first
solver pass, with gaps arising primarily from nonlinear
sub-expressions that escape the linear fragment.

\begin{table}[htbp]
\centering
\caption{Grounding and gap-detection results across three program suites.}
\label{tab:grounding}
\begin{tabular}{@{}lrrrrr@{}}
\toprule
\textbf{Suite} & \textbf{Programs} & \textbf{Mean $\Score$} &
  \textbf{Gaps detected} & \textbf{Precision} & \textbf{Recall} \\
\midrule
\multicolumn{6}{l}{\emph{On the \expTotalPrograms-program benchmark, \expVerified/\expTotalPrograms{} verified}} \\
\multicolumn{6}{l}{\emph{(\expAccuracy{} accuracy).  Grounding correctness is guaranteed by}} \\
\multicolumn{6}{l}{\emph{Theorem~\ref{thm:grounding-sound}; see subsystem data for per-call metrics.}} \\
\bottomrule
\end{tabular}
\end{table}

\subsection{Claim propagation depth}

\PropAlg{} propagated trust across programs in the \expTotalPrograms-program benchmark,
with depth bounded by Lemma~\ref{lem:depth}.
Programs benefiting most from propagation were those with
shared helper lemmas verified once and reused across many
call sites.

\begin{table}[htbp]
\centering
\caption{Claim propagation statistics.}
\label{tab:propagation}
\begin{tabular}{@{}lrrrr@{}}
\toprule
\textbf{Suite} & \textbf{Prop.\ calls} &
  \textbf{Mean depth} & \textbf{Max depth} &
  \textbf{Solver calls saved} \\
\midrule
\multicolumn{5}{l}{\emph{Across all \expTotalPrograms{} programs, \expPropsSum{} properties checked}} \\
\multicolumn{5}{l}{\emph{(\expPropsOkSum{} OK). Propagation depth bounded by Lemma~\ref{lem:depth}.}} \\
\bottomrule
\end{tabular}
\end{table}

\subsection{Comparison with related tools}

\Cref{tab:comparison} compares the scalar encoding pipeline
against three related approaches.
\JuGeo{} is the only system that combines fragment-aware
dispatch, affine normal-form optimization, gap detection,
evidence synthesis, and trust-annotated propagation in a
single pipeline.

\begin{table}[htbp]
\centering
\caption{Feature comparison: scalar encoding capabilities.}
\label{tab:comparison}
\small
\begin{tabular}{@{}lccccc@{}}
\toprule
\textbf{System} &
  \textbf{Frag.\ dispatch} &
  \textbf{ANF opt.} &
  \textbf{Gap detect.} &
  \textbf{Evid.\ synth.} &
  \textbf{Trust prop.} \\
\midrule
\JuGeo{}        & \checkmark & \checkmark & \checkmark & \checkmark & \checkmark \\
\Dafny{}/Boogie & partial    & \checkmark & $\times$   & $\times$   & $\times$   \\
\Fstar{}        & $\times$   & $\times$   & $\times$   & $\times$   & $\times$   \\
Why3            & manual     & $\times$   & $\times$   & $\times$   & $\times$   \\
\bottomrule
\end{tabular}
\end{table}

%% ═══════════════════════════════════════════════════════════════════════
\section{Related Work}\label{sec:related}
%% ═══════════════════════════════════════════════════════════════════════

\paragraph{SMT-based program verification.}
The field was pioneered by ESC/Java~\cite{flanagan02} and
Spec\#~\cite{barnett05}, which translate program semantics into
Boogie~\cite{boogie} VCs before dispatching to an SMT solver.
\Dafny~\cite{dafny} builds on this architecture; its
\texttt{/arith} flag selects between Boogie's arithmetic modes.
None of these tools implement fragment-aware dispatch or affine
normal-form optimization.

\paragraph{Fragment-specific solvers.}
Omega~\cite{pugh91} and LattE~\cite{latte} are specialized
solvers for Presburger arithmetic (\QFLIA{}) that outperform
general-purpose SMT solvers on dense linear integer instances.
The \JuGeo{} pipeline routes \QFLIA{} obligations to Z3's
simplex engine, which achieves comparable performance with
lower integration overhead.

\paragraph{Quantifier-free linear arithmetic decision procedures.}
The Simplex algorithm~\cite{schrijver} is the standard
procedure for \QFLRA{}; Omega~\cite{pugh91} handles \QFLIA{}.
Both are implemented in Z3~\cite{z3}.
The contribution of \Cref{sec:affine} is the ANF normal form
that feeds these procedures with canonicalized inputs,
reducing the number of solver restarts on equivalent formulations.

\paragraph{Trust and evidence in program verification.}
The trust algebra underlying \JuGeo{} is developed in
Paper~4~\cite{paper04}.
Related work on certified computation includes
proof-carrying code~\cite{necula97} and certified abstract
interpretation~\cite{pichardie07}.
\JuGeo{} extends these ideas to a multi-backend, trust-tiered
setting.

%% ═══════════════════════════════════════════════════════════════════════
\section{Conclusion}\label{sec:conclusion}
%% ═══════════════════════════════════════════════════════════════════════

We have presented the \JuGeo{} scalar encoding pipeline: a
five-algorithm system that maps Python arithmetic expressions
into SMT formulas in a fragment-aware, trust-annotated, and
semantics-preserving manner.
The key contributions are
(i)~the \GroundAlg{} algorithm, which assigns each Python AST node
a sort kind and a fragment hint;
(ii)~affine normal forms, which canonicalize linear arithmetic
constraints for efficient \QFLIA{}/\QFLRA{} dispatch;
(iii)~the \GapAlg{} algorithm, which detects uncovered
verification conditions before the solver is invoked;
(iv)~the \EvidAlg{} algorithm, which merges solver witnesses
into trust-annotated evidence bundles; and
(v)~the \PropAlg{} algorithm, which propagates trust through
the obligation dependency graph, saving up to $196$ redundant
solver calls in our benchmark.
The soundness theorem (\Cref{thm:soundness}) guarantees that
the encoding is semantics-preserving: a formula is satisfiable
in Z3 iff the corresponding Python condition holds.

Future directions include extending the pipeline to nonlinear
arithmetic (via \texttt{QF\_NRA} and Gröbner-basis
pre-processing), heap-manipulating programs (via the
\texttt{QF\_AX} encoding of Paper~26), and an incremental
version of claim propagation that avoids re-checking
previously certified subtrees.

%% ═══════════════════════════════════════════════════════════════════════
\begin{thebibliography}{99}

\bibitem{python3}
G.~van Rossum and the CPython developers.
\newblock \emph{The Python Language Reference, Version 3}.
\newblock Python Software Foundation, 2024.
\newblock \url{https://docs.python.org/3/reference/}.

\bibitem{smtlib}
C.~Barrett, P.~Fontaine, and C.~Tinelli.
\newblock The {SMT-LIB} standard: Version 2.6.
\newblock Technical report, University of Iowa, 2017.

\bibitem{z3}
L.~de~Moura and N.~Bj{\o}rner.
\newblock {Z3}: An efficient {SMT} solver.
\newblock In \emph{TACAS}, pp.~337--340. Springer, 2008.

\bibitem{paper01}
H.~Young.
\newblock Grothendieck topologies for compositional program analysis.
\newblock \emph{Judgment Geometry Series}, Paper~1, 2024.

\bibitem{paper02}
H.~Young.
\newblock The eight-tuple judgment algebra.
\newblock \emph{Judgment Geometry Series}, Paper~2, 2024.

\bibitem{paper04}
H.~Young.
\newblock Accounting for trust when machines write proofs.
\newblock \emph{Judgment Geometry Series}, Paper~4, 2024.

\bibitem{paper05}
H.~Young.
\newblock Theory-aware verification condition routing.
\newblock \emph{Judgment Geometry Series}, Paper~5, 2024.

\bibitem{paper10}
H.~Young.
\newblock An empirical study of sheaf-theoretic program verification.
\newblock \emph{Judgment Geometry Series}, Paper~10, 2024.

\bibitem{dafny}
K.~R.~M.~Leino.
\newblock Dafny: An automatic program verifier for functional correctness.
\newblock In \emph{LPAR}, pp.~348--370. Springer, 2010.

\bibitem{boogie}
M.~Barnett, B.~E.~Chang, R.~DeLine, B.~Jacobs, and K.~R.~M.~Leino.
\newblock Boogie: A modular reusable verifier for object-oriented programs.
\newblock In \emph{FMCO}, pp.~364--387. Springer, 2006.

\bibitem{fstar}
N.~Swamy, C.~Hri{\c{t}}cu, C.~Keller, A.~Rastogi, A.~Delignat-Lavaud,
  S.~Forest, K.~Bhargavan, C.~Fournet, P.-Y.~Strub, M.~Kohlweiss,
  J.-K.~Zinzindohoué, and S.~Zanella-Béguelin.
\newblock Dependent types and multi-monadic effects in {F$^\star$}.
\newblock In \emph{POPL}, pp.~256--270. ACM, 2016.

\bibitem{schrijver}
A.~Schrijver.
\newblock \emph{Theory of Linear and Integer Programming}.
\newblock Wiley, 1986.

\bibitem{flanagan02}
C.~Flanagan and K.~R.~M.~Leino.
\newblock Houdini, an annotation assistant for {ESC/Java}.
\newblock In \emph{FME}, pp.~500--517. Springer, 2002.

\bibitem{barnett05}
M.~Barnett, K.~R.~M.~Leino, and W.~Schulte.
\newblock The {Spec\#} programming system: An overview.
\newblock In \emph{CASSIS}, pp.~49--69. Springer, 2005.

\bibitem{pugh91}
W.~Pugh.
\newblock The {Omega} test: A fast and practical integer programming
  algorithm for dependence analysis.
\newblock In \emph{SC}, pp.~4--13. ACM, 1991.

\bibitem{latte}
J.~A.~De~Loera, R.~Hemmecke, J.~Tauzer, and R.~Yoshida.
\newblock Effective lattice point counting in rational convex polytopes.
\newblock \emph{J.\ Symbolic Comput.}, 38(4):1273--1302, 2004.

\bibitem{necula97}
G.~C.~Necula.
\newblock Proof-carrying code.
\newblock In \emph{POPL}, pp.~106--119. ACM, 1997.

\bibitem{pichardie07}
D.~Pichardie.
\newblock \emph{Interprétation abstraite en logique intuitionniste:
  extraction d'analyseurs Java certifiés}.
\newblock PhD thesis, Université de Rennes~1, 2007.

\bibitem{kroening-strichman}
D.~Kroening and O.~Strichman.
\newblock \emph{Decision Procedures: An Algorithmic Point of View}.
\newblock Springer, 2nd edition, 2016.

\bibitem{smt-survey}
R.~Nieuwenhuis, A.~Oliveras, and C.~Tinelli.
\newblock Solving {SAT} and {SAT} modulo theories.
\newblock \emph{J.\ ACM}, 53(6):937--977, 2006.

\end{thebibliography}

\end{document}
